\documentclass[11pt,a4paper]{article}
\usepackage[utf8]{inputenc}
\usepackage[T1]{fontenc}
\usepackage{lmodern}
\usepackage[spanish,es-tabla,es-noshorthands]{babel}
\usepackage[margin=1.4cm]{geometry}
\usepackage{tikz}
\usetikzlibrary{shapes.geometric, shapes.misc, arrows.meta, positioning, calc, fit, backgrounds, shadows.blur, babel}
\usepackage{xcolor}
\usepackage{booktabs}
\usepackage{tabularx}
\usepackage{amsmath,amssymb}
\usepackage{fancyhdr}
\usepackage{tcolorbox}

% Paleta de colores institucional UTI
\definecolor{utiNavy}{HTML}{0F2A4A}
\definecolor{utiBlue}{HTML}{1A5694}
\definecolor{utiLightBlue}{HTML}{EBF3FB}
\definecolor{criticalRed}{HTML}{DC2626}
\definecolor{criticalRedLight}{HTML}{FEF2F2}
\definecolor{slackGreen}{HTML}{16A34A}
\definecolor{slackGreenLight}{HTML}{F0FDF4}
\definecolor{borderSlate}{HTML}{94A3B8}

\pagestyle{fancy}
\fancyhf{}
\fancyhead[L]{\scriptsize\textbf{\color{utiNavy}UNIVERSIDAD TECNOLÓGICA INDOAMÉRICA} $\vert$ Ing. en Tecnologías de la Información}
\fancyhead[R]{\scriptsize\color{utiNavy}Gestión de Proyectos $\vert$ Método de la Ruta Crítica (CPM)}
\fancyfoot[L]{\scriptsize\color{gray}Autor: Cristofer Alexander Paez Caza (1754161709)}
\fancyfoot[R]{\thepage}
\renewcommand{\headrulewidth}{0.6pt}
\renewcommand{\footrulewidth}{0.4pt}

\begin{document}

% ==============================================================================
% ENCABEZADO FORMAL
% ==============================================================================
\begin{tcolorbox}[colback=utiNavy, colframe=utiNavy, arc=4pt, colupper=white, boxrule=0pt, left=10pt, right=10pt, top=8pt, bottom=8pt]
\begin{center}
    {\scriptsize\textbf{\color{utiLightBlue}FACULTAD DE INGENIERÍAS Y TECNOLOGÍAS DE LA INFORMACIÓN \quad $\bullet$ \quad UTI}}\\[3pt]
    {\large\textbf{\color{white}MODELADO Y EVALUACIÓN DE RUTA CRÍTICA (CPM / PERT)}}\\[2pt]
    {\small\textbf{\color{utiLightBlue}TRABAJO DE TITULACIÓN: CRONOGRAMA DE 16 SEMANAS}}\\[4pt]
    {\footnotesize\color{white}\textbf{Tema:} Actualización de un sitio web mediante diseño centrado en el usuario para mejorar la accesibilidad y la experiencia de personas con discapacidad}\\
    {\scriptsize\color{utiLightBlue}\textbf{Autor:} Cristofer Alexander Paez Caza \quad $\vert$ \quad \textbf{Tutora:} Ing. María Cristina Jiménez, Mg. \quad $\vert$ \quad \textbf{Período:} Octubre 2026 -- Enero 2027}
\end{center}
\end{tcolorbox}

\vspace{0.3cm}

% ==============================================================================
% 1. FUNDAMENTOS TEÓRICOS DE LA RED CPM
% ==============================================================================
\section*{\normalsize\color{utiNavy}1. Estructura del Nodo y Parámetros del Método de la Ruta Crítica (CPM)}

El método de la ruta crítica (CPM, por sus siglas en inglés) utiliza un diagrama de red de actividades en los nodos (AON). Cada actividad $i$ se modela mediante un bloque estructurado en 6 cuadrantes:

\begin{center}
\begin{tikzpicture}[every node/.style={font=\footnotesize}]
    \draw[line width=1.1pt, draw=utiNavy] (0,0) rectangle (6.6,1.8);
    \draw[line width=0.7pt, draw=utiNavy] (0,0.9) -- (6.6,0.9);
    \draw[line width=0.7pt, draw=utiNavy] (2.2,0) -- (2.2,1.8);
    \draw[line width=0.7pt, draw=utiNavy] (4.4,0) -- (4.4,1.8);
    
    % Fila superior (Ida)
    \node at (1.1, 1.35) {\textbf{IT}};
    \node at (3.3, 1.35) {\textbf{Duración ($D$)}};
    \node at (5.5, 1.35) {\textbf{FT}};
    
    % Fila inferior (Regreso)
    \node at (1.1, 0.45) {\textbf{IL}};
    \node at (3.3, 0.45) {\textbf{Holgura ($H$)}};
    \node at (5.5, 0.45) {\textbf{FL}};
    
    % Etiqueta
    \node[anchor=north, font=\bfseries\scriptsize, text=utiNavy] at (3.3, -0.1) {NOMBRE O CÓDIGO DE LA ACTIVIDAD};
\end{tikzpicture}
\end{center}

\vspace{-0.2cm}

\begin{itemize}\setlength\itemsep{1pt}
    \item \textbf{IT (Inicio Temprano / Early Start - ES):} El momento más temprano en que la actividad puede iniciar tras terminar sus predecesores inmediatos. $\text{IT} = \max(\text{FT}_{\text{predecesores}})$. Para la primera actividad, $\text{IT} = 0$.
    \item \textbf{FT (Fin Temprano / Early Finish - EF):} El momento más temprano en que la actividad puede culminar. $\text{FT} = \text{IT} + D$.
    \item \textbf{FL (Fin Tardío / Late Finish - LF):} El momento más tardío en que la actividad puede finalizar sin retrasar la fecha de culminación del proyecto. $\text{FL} = \min(\text{IL}_{\text{sucesores}})$. Para la última actividad, $\text{FL} = \text{FT}$.
    \item \textbf{IL (Inicio Tardío / Late Start - LS):} El momento más tardío en que la actividad puede iniciar sin retrasar el proyecto. $\text{IL} = \text{FL} - D$.
    \item \textbf{Holgura Total ($H$):} Margen de reprogramación o holgura de la actividad:
    \begin{equation}
        H = \text{IL} - \text{IT} = \text{FL} - \text{FT}
    \end{equation}
    \item \textbf{Ruta Crítica:} Cadena continua de actividades desde el inicio hasta el fin donde \textbf{la holgura es cero ($H = 0$)}. Cualquier retraso en estas actividades aplaza directamente la titulación.
\end{itemize}

% ==============================================================================
% 2. TABLA DE ACTIVIDADES DEL PROYECTO DE TITULACIÓN
% ==============================================================================
\section*{\normalsize\color{utiNavy}2. Matriz de Actividades y Dependencias del Proyecto (16 Semanas)}

\begin{center}
\footnotesize
\renewcommand{\arraystretch}{1.1}
\begin{tabularx}{\textwidth}{c X c c}
\toprule
\textbf{Cód.} & \textbf{Descripción de la Actividad en el Trabajo de Titulación} & \textbf{Predecesor} & \textbf{Dur. ($D$)} \\
\midrule
\textbf{A} & Diagnóstico técnico de laboratorio (Lighthouse/WAVE) y setup de Git & Ninguno & 2 sem \\
\textbf{B} & Subsanación y aprobación del Capítulo I (Contextualización y LOD/CONADIS) & A & 2 sem \\
\textbf{C} & Metodología DCU (ISO 9241-210) y caracterización de usuarios con discapacidad & B & 2 sem \\
\textbf{D} & Especificación formal de Requisitos (RF01--RF06 / RNF01--RNF05) y alcance & C & 1 sem \\
\textbf{E} & Prototipado interactivo accesible UI/UX en baja y alta fidelidad & D & 1 sem \\
\textbf{F} & Arquitectura modular frontend: Capa 1 (UI WAI-ARIA) y Capa 2 (Lógica Foco) & E & 2 sem \\
\textbf{G} & Desarrollo de Capa 3: Servicios RESTful y consumo de endpoints de datos & E & 1 sem \\
\textbf{H} & Optimización técnica y rendimiento: Core Web Vitals (TBT, LCP, WebP, Lazy) & F, G & 2 sem \\
\textbf{I} & Redacción preliminar de la memoria del Capítulo III (Implementación de software) & F & 1 sem \\
\textbf{J} & Re-auditorías técnicas automatizadas de control (Lighthouse $\ge 90$, WAVE = 0) & H & 1 sem \\
\textbf{K} & Pruebas empíricas con usuarios con discapacidad (NVDA/Teclado) y escala SUS & H & 2 sem \\
\textbf{L} & Capítulo IV (Resultados), Conclusiones, formato UTI y solicitud de lectores & I, J, K & 2 sem \\
\bottomrule
\end{tabularx}
\end{center}

\newpage

% ==============================================================================
% 3. CÁLCULO PASO A PASO (IDA, REGRESO Y HOLGURAS)
% ==============================================================================
\section*{\normalsize\color{utiNavy}3. Memoria de Cálculo Matemático Paso a Paso}

\subsection*{\small\color{utiBlue}A. Paso Hacia Adelante (Forward Pass) $\rightarrow$ Tiempos Tempranos ($\text{FT} = \text{IT} + D$)}
\begin{enumerate}\setlength\itemsep{2pt}\footnotesize
    \item \textbf{Actividad A:} $\text{IT}_A = 0 \implies \text{FT}_A = 0 + 2 = \mathbf{2}$.
    \item \textbf{Actividad B:} Viene de A: $\text{IT}_B = \text{FT}_A = 2 \implies \text{FT}_B = 2 + 2 = \mathbf{4}$.
    \item \textbf{Actividad C:} Viene de B: $\text{IT}_C = \text{FT}_B = 4 \implies \text{FT}_C = 4 + 2 = \mathbf{6}$.
    \item \textbf{Actividad D:} Viene de C: $\text{IT}_D = \text{FT}_C = 6 \implies \text{FT}_D = 6 + 1 = \mathbf{7}$.
    \item \textbf{Actividad E:} Viene de D: $\text{IT}_E = \text{FT}_D = 7 \implies \text{FT}_E = 7 + 1 = \mathbf{8}$.
    \item \textbf{Actividad F:} Viene de E: $\text{IT}_F = \text{FT}_E = 8 \implies \text{FT}_F = 8 + 2 = \mathbf{10}$.
    \item \textbf{Actividad G (En paralelo con F):} Viene de E: $\text{IT}_G = \text{FT}_E = 8 \implies \text{FT}_G = 8 + 1 = \mathbf{9}$.
    \item \textbf{Actividad H (Convergencia):} Predecesores F ($\text{FT}=10$) y G ($\text{FT}=9$). Aplicamos el \textbf{máximo}:\\
    $\text{IT}_H = \max(\text{FT}_F, \text{FT}_G) = \max(10, 9) = \mathbf{10} \implies \text{FT}_H = 10 + 2 = \mathbf{12}$.
    \item \textbf{Actividad I (En paralelo con H):} Viene de F: $\text{IT}_I = \text{FT}_F = 10 \implies \text{FT}_I = 10 + 1 = \mathbf{11}$.
    \item \textbf{Actividad J:} Viene de H: $\text{IT}_J = \text{FT}_H = 12 \implies \text{FT}_J = 12 + 1 = \mathbf{13}$.
    \item \textbf{Actividad K (En paralelo con J):} Viene de H: $\text{IT}_K = \text{FT}_H = 12 \implies \text{FT}_K = 12 + 2 = \mathbf{14}$.
    \item \textbf{Actividad L (Convergencia Final):} Predecesores I ($\text{FT}=11$), J ($\text{FT}=13$) y K ($\text{FT}=14$).\\
    $\text{IT}_L = \max(\text{FT}_I, \text{FT}_J, \text{FT}_K) = \max(11, 13, 14) = \mathbf{14} \implies \text{FT}_L = 14 + 2 = \mathbf{16}$.
\end{enumerate}
\vspace{-0.1cm}
\begin{tcolorbox}[colback=utiLightBlue, colframe=utiBlue, arc=3pt, left=8pt, right=8pt, top=4pt, bottom=4pt]
\footnotesize\textbf{Duración Total del Proyecto de Titulación:} $\mathbf{16\text{ semanas}}$ (Correspondiente a las 16 semanas calendario oficiales del semestre UTI).
\end{tcolorbox}

\subsection*{\small\color{utiBlue}B. Paso Hacia Atrás (Backward Pass) $\rightarrow$ Tiempos Tardíos ($\text{IL} = \text{FL} - D$)}
\begin{enumerate}\setlength\itemsep{2pt}\footnotesize
    \item \textbf{Actividad L (Última):} $\text{FL}_L = \text{FT}_L = 16 \implies \text{IL}_L = 16 - 2 = \mathbf{14}$.
    \item \textbf{Actividad K:} Sucesor L: $\text{FL}_K = \text{IL}_L = 14 \implies \text{IL}_K = 14 - 2 = \mathbf{12}$.
    \item \textbf{Actividad J:} Sucesor L: $\text{FL}_J = \text{IL}_L = 14 \implies \text{IL}_J = 14 - 1 = \mathbf{13}$.
    \item \textbf{Actividad I:} Sucesor L: $\text{FL}_I = \text{IL}_L = 14 \implies \text{IL}_I = 14 - 1 = \mathbf{13}$.
    \item \textbf{Actividad H (Divergencia de retorno):} Sucesores J ($\text{IL}=13$) y K ($\text{IL}=12$). Aplicamos el \textbf{mínimo}:\\
    $\text{FL}_H = \min(\text{IL}_J, \text{IL}_K) = \min(13, 12) = \mathbf{12} \implies \text{IL}_H = 12 - 2 = \mathbf{10}$.
    \item \textbf{Actividad G:} Sucesor H: $\text{FL}_G = \text{IL}_H = 10 \implies \text{IL}_G = 10 - 1 = \mathbf{9}$.
    \item \textbf{Actividad F (Divergencia de retorno):} Sucesores H ($\text{IL}=10$) e I ($\text{IL}=13$).\\
    $\text{FL}_F = \min(\text{IL}_H, \text{IL}_I) = \min(10, 13) = \mathbf{10} \implies \text{IL}_F = 10 - 2 = \mathbf{8}$.
    \item \textbf{Actividad E (Divergencia de retorno):} Sucesores F ($\text{IL}=8$) y G ($\text{IL}=9$).\\
    $\text{FL}_E = \min(\text{IL}_F, \text{IL}_G) = \min(8, 9) = \mathbf{8} \implies \text{IL}_E = 8 - 1 = \mathbf{7}$.
    \item \textbf{Actividad D:} Sucesor E: $\text{FL}_D = \text{IL}_E = 7 \implies \text{IL}_D = 7 - 1 = \mathbf{6}$.
    \item \textbf{Actividad C:} Sucesor D: $\text{FL}_C = \text{IL}_D = 6 \implies \text{IL}_C = 6 - 2 = \mathbf{4}$.
    \item \textbf{Actividad B:} Sucesor C: $\text{FL}_B = \text{IL}_C = 4 \implies \text{IL}_B = 4 - 2 = \mathbf{2}$.
    \item \textbf{Actividad A:} Sucesor B: $\text{FL}_A = \text{IL}_B = 2 \implies \text{IL}_A = 2 - 2 = \mathbf{0}$.
\end{enumerate}

% ==============================================================================
% 4. TABLA RESUMEN MAESTRA DE VALORES CPM
% ==============================================================================
\section*{\normalsize\color{utiNavy}4. Tabla Maestra de Resultados CPM}

\begin{center}
\small
\begin{tabularx}{\textwidth}{c X c c c c c c c}
\toprule
\textbf{Act.} & \textbf{Descripción Resumida} & \textbf{Dur.} & \textbf{IT} & \textbf{FT} & \textbf{IL} & \textbf{FL} & \textbf{Holgura ($H$)} & \textbf{¿Crítica?} \\
\midrule
\textbf{A} & Diagnóstico técnico de laboratorio y setup & 2 & 0 & 2 & 0 & 2 & $0 - 0 = \mathbf{0}$ & \textbf{\color{criticalRed}SÍ (Ruta Crítica)} \\
\textbf{B} & Subsanación y aprobación de Cap. I & 2 & 2 & 4 & 2 & 4 & $2 - 2 = \mathbf{0}$ & \textbf{\color{criticalRed}SÍ (Ruta Crítica)} \\
\textbf{C} & Metodología DCU y contexto de usuarios & 2 & 4 & 6 & 4 & 6 & $4 - 4 = \mathbf{0}$ & \textbf{\color{criticalRed}SÍ (Ruta Crítica)} \\
\textbf{D} & Especificación formal de Requisitos & 1 & 6 & 7 & 6 & 7 & $6 - 6 = \mathbf{0}$ & \textbf{\color{criticalRed}SÍ (Ruta Crítica)} \\
\textbf{E} & Prototipado interactivo UI/UX & 1 & 7 & 8 & 7 & 8 & $7 - 7 = \mathbf{0}$ & \textbf{\color{criticalRed}SÍ (Ruta Crítica)} \\
\textbf{F} & Arquitectura frontend: Capas 1 y 2 & 2 & 8 & 10 & 8 & 10 & $8 - 8 = \mathbf{0}$ & \textbf{\color{criticalRed}SÍ (Ruta Crítica)} \\
\textbf{G} & Capa 3: Servicios REST y datos & 1 & 8 & 9 & 9 & 10 & $9 - 8 = \mathbf{1}$ & \color{slackGreen}No (1 sem) \\
\textbf{H} & Optimización Core Web Vitals & 2 & 10 & 12 & 10 & 12 & $10 - 10 = \mathbf{0}$ & \textbf{\color{criticalRed}SÍ (Ruta Crítica)} \\
\textbf{I} & Redacción borrador Capítulo III & 1 & 10 & 11 & 13 & 14 & $13 - 10 = \mathbf{3}$ & \color{slackGreen}No (3 sem) \\
\textbf{J} & Re-auditorías técnicas (Lighthouse/WAVE) & 1 & 12 & 13 & 13 & 14 & $13 - 12 = \mathbf{1}$ & \color{slackGreen}No (1 sem) \\
\textbf{K} & Pruebas con usuarios PCD y escala SUS & 2 & 12 & 14 & 12 & 14 & $12 - 12 = \mathbf{0}$ & \textbf{\color{criticalRed}SÍ (Ruta Crítica)} \\
\textbf{L} & Capítulo IV, Formato y Defensa de Grado & 2 & 14 & 16 & 14 & 16 & $14 - 14 = \mathbf{0}$ & \textbf{\color{criticalRed}SÍ (Ruta Crítica)} \\
\bottomrule
\end{tabularx}
\end{center}

\newpage

% ==============================================================================
% 5. DIAGRAMA DE RED DEL PROYECTO (AON)
% ==============================================================================
\section*{\normalsize\color{utiNavy}5. Diagrama de Red del Proyecto (AON - Activity On Node)}

En el siguiente diagrama se presenta la red completa de actividades con sus 6 cuadrantes de tiempos calculados. Las flechas rojas dobles y nodos destacados representan la \textbf{Ruta Crítica}:

\begin{center}
\resizebox{\textwidth}{!}{
\begin{tikzpicture}[
    % Estilo del nodo estándar
    cpmNode/.style={
        rectangle, rounded corners=2pt, draw=utiBlue, line width=1pt, fill=white,
        inner sep=0pt, minimum width=2.4cm, minimum height=1.3cm
    },
    cpmCritical/.style={
        rectangle, rounded corners=2pt, draw=criticalRed, line width=1.5pt, fill=criticalRedLight,
        inner sep=0pt, minimum width=2.4cm, minimum height=1.3cm
    },
    % Flechas
    arrowNorm/.style={->, >=Stealth, line width=1pt, color=borderSlate},
    arrowCrit/.style={->, >=Stealth, line width=2.2pt, color=criticalRed}
]

% Macro para dibujar el bloque de 6 cuadrantes
% #1: nombre nodo, #2: coords, #3: dur, #4: IT, #5: FT, #6: IL, #7: H, #8: FL, #9: estilo
\newcommand{\cpmBox}[9]{
    \node[#9] (box_#1) at #2 {
        \begin{tabular}{c|c|c}
        \footnotesize #4 & \footnotesize\bfseries #3 & \footnotesize #5 \\
        \hline
        \footnotesize #6 & \footnotesize\bfseries #7 & \footnotesize #8
        \end{tabular}
    };
    \node[above=2pt of box_#1, font=\bfseries\scriptsize, text=utiNavy] {\textbf{Act. #1}};
}

% Nodos de la Red
% Fila Principal
\cpmBox{A}{(0, 6)}{2}{0}{2}{0}{0}{2}{cpmCritical}
\cpmBox{B}{(3.2, 6)}{2}{2}{4}{2}{0}{4}{cpmCritical}
\cpmBox{C}{(6.4, 6)}{2}{4}{6}{4}{0}{6}{cpmCritical}
\cpmBox{D}{(9.6, 6)}{1}{6}{7}{6}{0}{7}{cpmCritical}
\cpmBox{E}{(12.8, 6)}{1}{7}{8}{7}{0}{8}{cpmCritical}

% Ramificaciones en Desarrollo
\cpmBox{F}{(16.2, 7.5)}{2}{8}{10}{8}{0}{10}{cpmCritical}
\cpmBox{G}{(16.2, 4.5)}{1}{8}{9}{9}{1}{10}{cpmNode}

% Convergencia en Optimización
\cpmBox{H}{(19.6, 6)}{2}{10}{12}{10}{0}{12}{cpmCritical}

% Paralelo Redacción Cap. III
\cpmBox{I}{(19.6, 9.2)}{1}{10}{11}{13}{3}{14}{cpmNode}

% Ramificaciones en Pruebas
\cpmBox{J}{(23.0, 7.5)}{1}{12}{13}{13}{1}{14}{cpmNode}
\cpmBox{K}{(23.0, 4.5)}{2}{12}{14}{12}{0}{14}{cpmCritical}

% Cierre y Defensa
\cpmBox{L}{(26.4, 6)}{2}{14}{16}{14}{0}{16}{cpmCritical}

% Conexiones Críticas
\draw[arrowCrit] (box_A) -- (box_B);
\draw[arrowCrit] (box_B) -- (box_C);
\draw[arrowCrit] (box_C) -- (box_D);
\draw[arrowCrit] (box_D) -- (box_E);

% Bifurcación E -> F y G
\draw[arrowCrit] (box_E.east) -- ++(0.4, 0) |- (box_F.west);
\draw[arrowNorm] (box_E.east) -- ++(0.4, 0) |- (box_G.west);

% Convergencia F, G -> H
\draw[arrowCrit] (box_F.east) -- ++(0.4, 0) |- (box_H.west);
\draw[arrowNorm] (box_G.east) -- ++(0.4, 0) |- (box_H.west);

% Salida F -> I (Redacción)
\draw[arrowNorm] (box_F.north) |- (box_I.west);

% Bifurcación H -> J y K
\draw[arrowNorm] (box_H.east) -- ++(0.4, 0) |- (box_J.west);
\draw[arrowCrit] (box_H.east) -- ++(0.4, 0) |- (box_K.west);

% Convergencia I, J, K -> L
\draw[arrowNorm] (box_I.east) -| ($(box_L.west)+(-0.4, 0)$) -- (box_L.west);
\draw[arrowNorm] (box_J.east) -- ++(0.4, 0) |- (box_L.west);
\draw[arrowCrit] (box_K.east) -- ++(0.4, 0) |- (box_L.west);

\end{tikzpicture}
}
\end{center}

\vspace{0.4cm}

% ==============================================================================
% 6. CONCLUSIONES PARA RESPONDER EN EL EXAMEN
% ==============================================================================
\section*{\normalsize\color{utiNavy}6. Respuestas Clave de Examen}

\begin{tcolorbox}[colback=white, colframe=utiNavy, arc=4pt, title={\textbf{\color{white}SÍNTESIS DE RESULTADOS OPERATIVOS}}, fonttitle=\bfseries\small]
\begin{enumerate}\setlength\itemsep{3pt}\small
    \item \textbf{Duración Total del Proyecto:} $\mathbf{16\text{ semanas}}$ (4 meses calendario).
    \item \textbf{Ruta Crítica Identificada:}
    $$\mathbf{A \longrightarrow B \longrightarrow C \longrightarrow D \longrightarrow E \longrightarrow F \longrightarrow H \longrightarrow K \longrightarrow L}$$
    \item \textbf{Actividades Críticas:} A, B, C, D, E, F, H, K, L. Su holgura es estrictamente cero ($H = 0$). Cualquier desviación en ellas aplaza la entrega o defensa de la tesis.
    \item \textbf{Actividades No Críticas con Margen de Holgura:}
    \begin{itemize}\footnotesize
        \item \textbf{Actividad G (Servicios REST):} Holgura de $\mathbf{1\text{ semana}}$. Puede durar hasta 2 semanas sin retrasar la integración de optimización (H).
        \item \textbf{Actividad I (Borrador Cap. III):} Holgura de $\mathbf{3\text{ semanas}}$. Al redactarse en paralelo con el software, no compromete la fecha final hasta la semana 14.
        \item \textbf{Actividad J (Re-auditorías automatizadas):} Holgura de $\mathbf{1\text{ semana}}$. Puede ejecutarse de forma flexible mientras concluyen las pruebas de usuarios (K).
    \end{itemize}
\end{enumerate}
\end{tcolorbox}

\end{document}
